\documentclass[a4paper,12pt]{article}

\usepackage[utf8]{inputenc}
\usepackage[T1]{fontenc}
\usepackage[table]{xcolor}
\usepackage{geometry}
\usepackage{longtable}
\usepackage{enumitem}
\usepackage{hyperref}
\usepackage{listings}
\usepackage{titlesec}
\usepackage{booktabs}
\usepackage{array}
\usepackage{amsmath}

\geometry{margin=2.5cm}

\definecolor{criticalcolor}{HTML}{DC3545}
\definecolor{highcolor}{HTML}{FD7E14}
\definecolor{mediumcolor}{HTML}{FFC107}
\definecolor{lowcolor}{HTML}{28A745}
\definecolor{resolvedcolor}{HTML}{6C757D}
\definecolor{newcolor}{HTML}{007BFF}
\definecolor{scorecolor}{HTML}{17A2B8}

\newcommand{\severitybox}[2]{%
  \colorbox{#1}{\makebox[2.5cm][c]{\textcolor{white}{\textbf{#2}}}}%
}

\newcommand{\scorebox}[2]{%
  \colorbox{#1}{\makebox[3.0cm][c]{\textcolor{white}{\textbf{#2}}}}%
}

\lstset{
  basicstyle=\small\ttfamily,
  breaklines=true,
  frame=single,
  numbers=none,
  tabsize=2,
  backgroundcolor=\color{gray!10},
  extendedchars=true,
  inputencoding=utf8,
  literate={—}{---}1 {·}{$\cdot$}1 {≈}{$\approx$}1 {×}{$\times$}1 {→}{$\rightarrow$}1
}

\title{\textbf{Performance Report \#13}\\[4pt]
\large Vulkan Engine — Deep Re-Audit \& New CPU/GPU
Performance Analysis (Rendering-Pipeline Headroom)}
\author{Henrique Rocha}
\date{\today}

\begin{document}
\maketitle

\begin{abstract}
Reports \#10 and \#11 raised 13 issues (GPU shader tree + CPU$\leftrightarrow$GPU
sync). A code re-audit confirms 10 are fully resolved, 1 is mitigated (S1 — the
per-chunk serialization is gone but a single in-flight upload slot remains), and
2 are still open (GS1, NL2, both minor ALU redundancies). With the synchronous
stall surface removed, the engine is now genuinely pipelined; the remaining
performance \emph{headroom} is no longer in host stalls but in the \emph{rendering
pipeline itself}. This report therefore shifts focus: it documents a \emph{new}
set of issues found by tracing the steady-state frame — the dominant one being
that the full scene geometry is re-rendered six times per frame for cascaded
EVSM shadows, plus a 128-bit shadow-moment format and a per-tessellated-vertex
wave/Depth fetch in the water stage. Each new finding includes code evidence, a
quantified rationale, and a recommendation. The application is re-scored at
\textbf{8.4/10} with a clear path to \textbf{9.5/10} by addressing the shadow and
water-stage costs. All recommendations remain compatible with \texttt{AGENTS.md}.
\end{abstract}

\tableofcontents
\newpage

\section{Brief Status of Prior Findings (\#10 + \#11)}

The synchronous-stall class of problem is solved. Summary only (full evidence in
the previous revision):

\begin{center}
\begin{tabular}{cllc}
\toprule
\textbf{ID} & \textbf{Issue} & \textbf{Orig. severity} & \textbf{Status} \\
\midrule
C1  & veg. instance-gen serial workgroup    & High   & \severitybox{resolvedcolor}{RESOLVED} \\
C2  & per-invocation plane normalization     & Medium & \severitybox{resolvedcolor}{RESOLVED} \\
NL1 & perlin2d cos/sin hash                  & Medium & \severitybox{resolvedcolor}{RESOLVED} \\
T1  & triplanar weights 3×/vertex           & Medium & \severitybox{resolvedcolor}{RESOLVED} \\
GS1 & Fibonacci view table per primitive    & Medium & \severitybox{mediumcolor}{OPEN} \\
NL2 & sky sin star hash                      & Low    & \severitybox{lowcolor}{OPEN} \\
\midrule
S1  & per-chunk upload fence serialization   & High   & \severitybox{resolvedcolor}{MITIGATED} \\
S2  & per-frame readback + queue idle        & High   & \severitybox{resolvedcolor}{RESOLVED} \\
S3  & staging buffer per chunk               & Medium & \severitybox{resolvedcolor}{RESOLVED} \\
S4  & fence-destroy queue idle               & Medium & \severitybox{resolvedcolor}{RESOLVED} \\
S5  & geometryQueue unused                   & Medium & \severitybox{resolvedcolor}{RESOLVED} \\
S6  & blocking GPU veg. gen                  & Low    & \severitybox{resolvedcolor}{RESOLVED} \\
S7  & water reset deviceWaitIdle             & Low    & \severitybox{resolvedcolor}{RESOLVED} \\
\bottomrule
\end{tabular}
\end{center}

The key takeaway: the engine no longer blocks the host. The 8-submit,
3-frames-in-flight model is intact and uploads overlap rendering. \emph{New}
analysis below targets where the GPU now spends its time.

\newpage

\section{New Issue Register — Rendering-Pipeline Headroom}

\subsection{GPU / Geometry Throughput}

\begin{enumerate}[leftmargin=*]

\item[\textbf{[N1]}] \textbf{Full Scene Geometry Re-Rendered 6×/Frame for Cascaded EVSM Shadows}
  \hfill \severitybox{highcolor}{HIGH}

  \begin{tabular}{ll}
    \textbf{File:} & \texttt{vulkan/renderer/ShadowRenderer.cpp} \\
    \textbf{Line:} & 551 (cascade loop), 538 (\texttt{vkCmdDrawIndexed} per mesh) \\
    \textbf{Category:} & Vertex/geometry throughput / redundant shading \\
  \end{tabular}

  \textbf{Description:}
  \texttt{render()} iterates \texttt{SHADOW\_CASCADE\_COUNT = 3}
  (\texttt{UniformObject.hpp:8}) cascades. For each cascade it performs a depth
  pass and an EVSM color pass, each drawing the full shadow-casting geometry:

\begin{lstlisting}
// ShadowRenderer.cpp:551
for (int c = 0; c < SHADOW_CASCADE_COUNT; ++c) {
    renderCascade(app, commandBuffer, c, ...);   // depth + EVSM color
    blurCascade(app, commandBuffer, c, ...);     // 2 full-screen blur passes
}
// renderCascade issues, per mesh:
vkCmdDrawIndexed(commandBuffer, vbo.indexCount, 1, 0, 0, 0); // :538
\end{lstlisting}

  Cascade sizes are \texttt{\{2048, 1024, 512\}}
  (\texttt{ShadowRenderer.cpp:20}), so the visible cost is dominated by cascade~0
  (2048×2048) but the vertex/index/setup work is essentially the \emph{same mesh
  list} across all three. Net: the scene's entire draw list is issued
  \textbf{6 times per frame} (3 cascades × 2 passes), plus 6 separable blur
  passes. This is by far the largest single consumer of GPU vertex and geometry
  throughput in the frame and the primary reason frame GPU time scales with
  scene size rather than with what is visible to the camera.

  \textbf{Recommendation:}
  Decouple shadow update rate from the render rate. The shadow map only needs to
  change when the light direction or the camera (cascade-bounding) transform
  changes meaningfully. Maintain a dirty flag: when neither moved beyond a
  threshold, skip \texttt{renderCascade}/\texttt{blurCascade} and reuse the last
  shadow textures (optionally with stabilized cascade splits to avoid swimming).
  For static or slow-moving cameras this drops the 6 geometry passes to 0 on most
  frames. A lighter variant: render only cascade~0 at full rate and update
  cascades~1–2 every $N$ frames. Identical output when the light is static; a
  sub-frame-latency shadow update when it moves.

\item[\textbf{[N2]}] \textbf{Water Wave Displacement Paid at Tessellation-Eval Density with a Per-Vertex Scene-Depth Fetch}
  \hfill \severitybox{mediumcolor}{MEDIUM}

  \begin{tabular}{ll}
    \textbf{File:} & \texttt{shaders/water.tese} \\
    \textbf{Line:} & 102--118 (depth fetch), 119--152 (FBM-gradient) \\
    \textbf{Category:} & Tessellation-eval ALU / dependent texture fetch \\
  \end{tabular}

  \textbf{Description:}
  The water surface displacement is computed in the tessellation-evaluation
  stage, i.e. once per \emph{generated} vertex (the densest vertex multiplicity in
  the pipeline). Two costs compound there:

\begin{lstlisting}
// water.tese:102 -- depth-dependent attenuation samples the SCENE depth
// texture per generated vertex, inside the evaluation stage:
float solidDepthRaw = texture(sceneDepthTex, screenUV).r;   // dependent TFetch
...
// water.tese:119 -- 4D Perlin FBM *gradient* (octaves of perlin4d):
vec4 wave = waterWaveSample(xyz, animTime, noiseScale,
                            noiseOctaves, noisePersistence,
                            noiseLacunarity, bumpAmp, waveScale);
\end{lstlisting}

  A texture fetch in the tessellation-evaluation stage has poor cache coherence
  (vertices are not screen-contiguous) and stalls the vertex pipeline; the
  \texttt{noiseOctaves}-deep 4D Perlin FBM gradient is then paid at the same
  highest vertex count. Both scale with tessellation factor, not with visible
  pixels.

  \textbf{Recommendation:}
  (a) Move the depth-dependent attenuation out of the evaluation stage: compute
  the shallow-water suppression in the water \texttt{.frag} shader (per pixel,
  where the scene-depth texture is coherent) and feed a small per-vertex factor
  only if needed. (b) Cap \texttt{noiseOctaves} (e.g. 3) and/or evaluate the low-
  frequency displacement once per patch corner in \texttt{water.vert} and let the
  tessellator interpolate it — acceptable because wave displacement is low-
  spatial-frequency. This removes the per-generated-vertex texture fetch and
  cuts the FBM cost to the control-point rate.

\item[\textbf{[N3]}] \textbf{EVSM Moments Stored as 128-bit RGBA32F for Every Cascade}
  \hfill \severitybox{mediumcolor}{MEDIUM}

  \begin{tabular}{ll}
    \textbf{File:} & \texttt{vulkan/renderer/ShadowRenderer.cpp} \\
    \textbf{Line:} & 17 (format), 70 (color image), 229 (blur temp) \\
    \textbf{Category:} & Memory bandwidth / shadow resolution \\
  \end{tabular}

  \textbf{Description:}
  The EVSM color target is \texttt{VK\_FORMAT\_R32G32B32A32\_SFLOAT}
  (RGBA32F):

\begin{lstlisting}
static constexpr VkFormat EVSM_FORMAT = VK_FORMAT_R32G32B32A32_SFLOAT; // :17
// cascade color image (:70) and the separable-blur ping-pong temp (:229)
// are both RGBA32F.
\end{lstlisting}

  Standard EVSM needs only \textbf{two} moments (EVSM2, R32G32), which is
  algebraically sufficient for the positive/negative exponential warping used by
  EVSM. Storing 4 channels at 128 bits/texel across
  $2048^2 + 1024^2 + 512^2 \approx 5.5$ M texels means $\sim$88 MB of 128-bit
  render+sample+blur bandwidth per frame is spent on shadow moments, half of it
  redundant. The blur ping-pong temp doubles the write traffic.

  \textbf{Recommendation:}
  Switch \texttt{EVSM\_FORMAT} to \texttt{VK\_FORMAT\_R32G32\_SFLOAT}
  (EVSM2) — the canonical, bandwidth-optimal EVSM variant. This halves all
  shadow-moment bandwidth (render, blur read/write, and sampling in the lighting
  shader) with no visible quality regression for typical scenes. The lighting
  shader's moment reconstruction must be updated from 4- to 2-moment accordingly.

\subsection{GPU / Fragment \& Minor}

\item[\textbf{[N4]}] \textbf{Residual Shader ALU Redundancies (carried from \#10)}
  \hfill \severitybox{mediumcolor}{MEDIUM}

  \begin{tabular}{ll}
    \textbf{File:} & \texttt{shaders/impostors.geom} (\textbf{GS1}), \texttt{shaders/sky.frag} (\textbf{NL2}) \\
    \textbf{Category:} & Geometry/fragment ALU \\
  \end{tabular}

  \textbf{Description:}
  Still open from \#10: the 20-entry Fibonacci view-selection table is rebuilt
  per emitted impostor primitive (\texttt{impostors.geom:67--95}), and the sky
  star field still uses \texttt{fract(sin(dot(...)))} (\texttt{sky.frag:59}). Both
  are pure ALU waste in hot-ish paths (impostor primitive emission scales with
  vegetation billboard count; the sky hash is per sky pixel).

  \textbf{Recommendation:}
  Hoist the Fibonacci directions to a compile-time \texttt{const vec3 fibDirs[20]}
  (or UBO) and replace the sky hash with the integer bit-mix already used in
  \texttt{perlin.glsl}. Small, safe, and finishes the \#09 hash migration.

\subsection{CPU / Per-Frame Uploads}

\item[\textbf{[N5]}] \textbf{Shadow \& Static UBOs Fully Re-Uploaded Every Frame}
  \hfill \severitybox{lowcolor}{LOW}

  \begin{tabular}{ll}
    \textbf{File:} & \texttt{vulkan/renderer/SceneRenderer.cpp} \\
    \textbf{Line:} & 253 (per-cascade UniformObject), 346 (static UBO) \\
    \textbf{Category:} & Avoidable host→device copy \\
  \end{tabular}

  \textbf{Description:}
  Every frame the renderer \texttt{memcpy}s a full \texttt{UniformObject} for
  each of the 3 cascades (line 253) and a full static UBO (line 346) into staging
  and copies them to the GPU, regardless of whether the camera/light transform
  changed:

\begin{lstlisting}
memcpy(uboStagingBuffers[frameIdx].map(stagingOff), &shadowUBO, sizeof(UniformObject)); // :253
memcpy(uboStagingBuffers[frameIdx].map(restoreOff), &uboStatic, sizeof(UniformObject)); // :346
\end{lstlisting}

  These writes are small, but they are unconditionally issued every frame even
  when the view-projection and light matrix are byte-identical to the previous
  frame (e.g. a paused/idle scene).

  \textbf{Recommendation:}
  Track a dirty flag on the view-projection and light transform; only re-upload
  the cascade/static UBOs when they actually change. Trivially avoids the copy
  and its barrier on idle frames, and pairs naturally with the \textbf{[N1]}
  shadow dirty flag.

\end{enumerate}

\newpage

\section{Deep Analysis: Where the GPU Spends Its Time Now}

With the host no longer stalling, the frame GPU budget is dominated by
\emph{geometry re-issue}, not by shader ALU:

\begin{center}
\begin{tabular}{lp{7cm}c}
\toprule
\textbf{Phase} & \textbf{Work} & \textbf{Relative cost} \\
\midrule
Cascaded shadows (N1+N3) & 3 cascades × (depth + EVSM color) full-scene draws
  + 2 blur passes each, RGBA32F & \textbf{Very high} \\
Scene opaque + water (N2) & forward scene, tessellated water eval/disp & High \\
Vegetation indirect        & GPU-culled indirect draws (healthy) & Medium \\
Impostors (N4/GS1)         & geometry-shader billboards, per-primitive table & Medium \\
Sky + post (N4/NL2)        & full-screen + sky frag & Low--Medium \\
\bottomrule
\end{tabular}
\end{center}

\noindent The decisive observation is that \textbf{[N1] alone accounts for more
draw calls than the rest of the frame combined} (6 full-scene geometry passes vs.
1 for the opaque scene). Because the shadow-casting mesh list is the same list
re-issued per cascade, and because cascade~0 is full-resolution, the marginal
cost of adding geometry to the scene is paid \emph{sixfold} every frame. Any
optimization that reduces shadow re-issuance (update-rate decoupling, smaller
cascade~0, EVSM2 bandwidth) compounds across the whole scene.

The water stage (\textbf{[N2]}) is the second area of interest: it is the only
place that issues a dependent texture fetch in a vertex/tessellation stage, which
is the worst stage for texture-cache behavior. Moving that fetch to the fragment
shader and lowering the FBM rate removes a stall-prone path.

\section{Coupled View}

These new issues are \emph{pure GPU-side headroom} — they do not stall the CPU.
The pipeline remains correctly pipelined (S1 residual aside). This is a much
healthier position than the pre-fix engine: the host is free to feed the GPU, and
the only question is how much geometry/shading the GPU must repeat. The
recommended changes are therefore throughput/quality trade-offs rather than
correctness fixes, and can be applied incrementally and measured with a GPU
timeline/profiler.

\newpage

\section{Application Scorecard (Re-Scored)}

Re-scored on the current tree, with the new rendering-pipeline headroom factored
in (the old sync issues are resolved, but the shadow/water costs are real):

\begin{center}
\begin{tabular}{lp{6.5cm}c}
\toprule
\textbf{Dimension} & \textbf{Basis} & \textbf{Score} \\
\midrule
Architecture \& Modernity & Dynamic rendering, Synchronization2, timeline
  semaphores, descriptor indexing, 3-frame in-flight, StagingRingBuffer & 9 / 10 \\
\midrule
GPU Shader Efficiency & C1/C2/NL1/T1 fixed; N1/N2/N3 headroom; GS1/NL2 open & 8 / 10 \\
\midrule
CPU$\leftrightarrow$GPU Sync & S2/S3/S4/S5/S6/S7 fixed; S1 single-slot residual & 8 / 10 \\
\midrule
CPU Compute Pipeline & Parallel octree meshing, no host stalls & 9 / 10 \\
\midrule
Memory \& Allocation Hygiene & Ring staging, fence free-list, no per-chunk alloc & 9 / 10 \\
\bottomrule
\end{tabular}
\end{center}

\subsection*{Weighted Overall Score}

\begin{center}
\begin{tabular}{lcc}
\toprule
\textbf{Dimension} & \textbf{Weight} & \textbf{Contribution} \\
\midrule
Architecture \& Modernity & 0.15 & 1.35 \\
GPU Shader Efficiency      & 0.25 & 2.00 \\
CPU$\leftrightarrow$GPU Sync & 0.30 & 2.40 \\
CPU Compute Pipeline      & 0.15 & 1.35 \\
Memory \& Allocation Hygiene & 0.15 & 1.35 \\
\midrule
\textbf{Overall} & \textbf{1.00} & \textbf{8.45 / 10} \\
\bottomrule
\end{tabular}
\end{center}

\scorebox{scorecolor}{OVERALL: 8.45 / 10}

\subsection*{Interpretation}

The application holds a strong \textbf{8.45/10 (A-)}. The synchronous-stall debt
from \#10/\#11 is paid; the engine is genuinely pipelined. The remaining gap to a
\textbf{9.5/10} is now exclusively in rendering-pipeline headroom, dominated by
\textbf{[N1]} (6× full-scene shadow re-issue) and amplified by \textbf{[N3]}
(128-bit shadow moments) and \textbf{[N2]} (per-vertex water depth fetch). These
are throughput optimizations with clean, low-risk implementations, not stall
fixes. Closing \textbf{[N1]}+\textbf{[N3]} would very likely lift the GPU Shader
Efficiency score to 9–10 and the overall into the 9s.

\newpage

\section{Roadmap}

\begin{center}
\begin{tabular}{cll}
\toprule
\textbf{Priority} & \textbf{Item} & \textbf{Expected impact} \\
\midrule
P0 & \textbf{[N1]} shadow update-rate decoupling / dirty flag & Removes up to 6 geom passes/frame \\
P1 & \textbf{[N3]} EVSM2 (R32G32) shadow-moment format        & Halves shadow bandwidth \\
P1 & \textbf{[N2]} move water depth fetch to frag; cap FBM    & Removes tess-eval TFetch stall \\
P2 & \textbf{[N4]} GS1 fib table + NL2 sky hash               & Finishes ALU cleanup \\
P2 & \textbf{[N5]} dirty-flag UBO re-uploads                  & Avoids idle-frame copies \\
P3 & \textbf{[S1]} upload slot ring                           & Fully async uploads \\
\bottomrule
\end{tabular}
\end{center}

\noindent All items preserve the invariants mandated by \texttt{AGENTS.md}:
Synchronization2, dynamic rendering, timeline semaphores, descriptor indexing,
feature detection, and the documented RADV same-queue fallbacks. The roadmap
prioritizes the two changes that compound across the entire scene (\textbf{[N1]},
\textbf{[N3]}) before the localized ALU cleanups.

\end{document}
